\documentclass[11pt]{article}
\usepackage[margin=1in]{geometry}
\usepackage{amsmath,amssymb,amsthm,booktabs}
\usepackage[hidelinks]{hyperref}
\usepackage{xurl}
\setlength{\emergencystretch}{2em}
\newtheorem{theorem}{Theorem}
\newtheorem{lemma}[theorem]{Lemma}
\newtheorem{corollary}[theorem]{Corollary}
\newcommand{\F}{\mathbb F}
\newcommand{\mr}{\operatorname{mr}}
\newcommand{\rk}{\operatorname{rank}}
\title{Ternary affine quotients with meridian class rank four}
\author{Exact finite-group proof and knot certificates}
\date{17 November 2026}
\begin{document}
\maketitle

A finite affine group of order $122472$ has ordinary generating rank
two, while its distinguished involution class requires four generators.
Twenty-four retained knot groups surject onto this group with all
meridians in that class. This gives a direct finite linear-group proof
of their already established equalities $\mr=b=4$.

The proof uses explicit matrices and kernel words. It requires no
branched-cover rank inequality or identification of a cover group.
These are the same twenty-four named equalities as in the retained
seven-sheeted-cover result; the number of named examples does not
increase. Worldwide priority, formal verification and external
specialist review are not claimed.

\section{A fixed-space bound}

For a finite-dimensional representation $\rho:H\to\mathrm{GL}(V)$,
write $V^H$ for its common fixed space.

\begin{lemma}\label{lem:fixed}
Suppose $H$ is generated by $k$ elements whose displacement matrices
$\rho(h_i)-I$ have rank at most $\delta$. Then
\[
                       \dim V-\dim V^H\le k\delta.
\]
Consequently a meridional representation of a knot group whose
meridians have displacement rank $\delta>0$ gives
\[
              \mr(K)\ge
              \left\lceil\frac{\dim V-\dim V^{\rho(\pi_1)}}{\delta}
              \right\rceil.
\]
\end{lemma}
\begin{proof}
The common fixed space is the intersection of the kernels of the
$k$ displacement maps. The codimension of an intersection is at most
the sum of the codimensions. Meridian images have the same
displacement rank by conjugacy. This proves both assertions.
\end{proof}

This bound applies to any field and does not require reflections,
a nondegenerate form or a faithful representation.

\section{The finite group and its involution class}

Let $G=\mathrm{GL}_3(\F_2)$ act on its seven nonzero vectors, numbered
$0,\ldots,6$ by binary values $1,\ldots,7$.
Let
\[
 W=\{(w_0,\ldots,w_6)\in\F_3^7:\textstyle\sum_iw_i=0\},
 \qquad Q=W\rtimes G.
\]
For a permutation $p$, let $P_pe_i=e_{p(i)}$. Multiplication is
\[
             (p,v)(q,w)=(pq,v+P_pw).
\]
Use the homogeneous action on $V=W\oplus\F_3$, of dimension seven:
\[
                   (p,v)(w,s)=(P_pw+sv,s).
\]
A convenient basis for $W$ is $f_j=e_j-e_6$, $0\le j<6$.
In this basis the homogeneous matrix is explicitly
\begin{equation}\label{eq:matrix}
 L(p,v)=
 \begin{pmatrix}
  (\mathbf1_{p(j)=i}-\mathbf1_{p(6)=i})_{0\le i,j<6}
       &(v_i)_{0\le i<6}\\
  0&1
 \end{pmatrix}.
\end{equation}
All entries are taken modulo three.

Let $C$ be the conjugacy class in $Q$ of $(p,0)$, where $p$ is any
nonidentity involution of $G$. This does not depend on the chosen
involution.

\begin{theorem}\label{thm:group}
The group $Q$ has order $3^6\cdot168=122472$. Its class $C$ has
$189$ elements and generating rank exactly four.
Every element of $C$ acts on $V$ with displacement rank two, and
$V^Q=0$. The ordinary generating rank of $Q$ is exactly two.
\end{theorem}
\begin{proof}
The group $G$ has $(8-1)(8-2)(8-4)=168$ elements and $W$ has dimension
six, proving the order formula.

A nonidentity involution $A$ in $\mathrm{GL}_3(\F_2)$ has
$(A-I)^2=0$ and $\rk(A-I)=1$. Thus $A-I=u\phi$, with $u\ne0$,
$\phi\ne0$ and $\phi(u)=0$. Adapted bases show that all such
involutions are conjugate. There are seven choices for $u$ and
three choices for a nonzero $\phi$ annihilating it, hence twenty-one
involutions. In the seven-point action each has three fixed points
and two transpositions.

Over $\F_3$, the negative eigenspace of such a permutation has
dimension two and lies in $W$. Since $7$ is nonzero modulo three,
the permutation module splits as the constant line and $W$.
Consequently $\rk(P_p-I)|_W=2$.
The affine involutions above $p$ are precisely
\[
                       (p,v),\qquad (I+P_p)v=0.
\]
Every such $v$ is obtained from a translation conjugate:
$(I-P_p)(v/2)=v$. Thus $C$ consists of these nine lifts of each of
the twenty-one involutions, so $|C|=21\cdot9=189$.
Conjugacy preserves the displacement rank of $L(p,0)$, which is two.

A vector $(w,s)$ fixed by all translations has $s=0$.
The only $G$-fixed vectors in $\F_3^7$ are constants, by transitivity;
a constant vector in $W$ is zero because $7\ne0$ in $\F_3$.
Therefore $V^Q=0$. Lemma~\ref{lem:fixed} shows that fewer than four
elements of $C$ cannot generate $Q$.

Here is an explicit generating quartet attaining four.
In cycle notation on $\{0,\ldots,6\}$ set
\[
 p=(0\,4)(1\,5),\qquad q=(0\,6)(2\,4),\qquad r=(1\,2)(3\,4),
 \qquad v=e_1-e_5.
\]
Take
\[
 x_0=(p,v),\qquad x_1=(p,0),\qquad
 x_2=(q,0),\qquad x_3=(r,0).
\]
They all belong to $C$. The three binary matrices for $p,q,r$
have row masks $(1,2,7),(1,3,5),(7,2,4)$ respectively;
the exact finite enumeration verifies that they generate all $168$
elements of $G$.
The following words in the involutions $x_i$ are pure translations:
\[
\begin{array}{c|c}
\text{word, written as a sequence of subscripts}&\text{translation}\\\hline
10&e_5-e_1\\
3103&e_5-e_2\\
231032&e_5-e_4\\
02310320&e_1-e_0\\
32310323&e_5-e_3\\
2023103202&e_1-e_6 .
\end{array}
\]
Their translations are independent and span $W$.
Direct word multiplication verifies every row.
The quartet therefore contains the full translation group and
surjects onto $G$, hence generates $Q$.

Finally, the two words $x_0x_2$ and $x_1x_3$ generate $Q$.
The independent finite control enumerates their affine closure on the
$168\cdot729$ possible pairs and obtains exactly $122472$ elements.
This is an exhaustive finite group calculation, retained in the
checker. Since $Q$ has a nonabelian quotient $G$, it is not cyclic.
Its ordinary generating rank is therefore two.
\end{proof}

\section{Twenty-four knot-group epimorphisms}

Put
\[
\begin{split}
\mathcal K=\{&
13a_{373},13a_{375},13a_{391},13a_{392},
13a_{3264},13a_{3266},13a_{3275},13a_{3276},\\
&13n_{1005},13n_{1006},13n_{1007},13n_{1011},13n_{1012},13n_{1013},
13n_{1043},13n_{1044},\\
&13n_{1045},13n_{1046},13n_{1047},13n_{1048},
13n_{3842},13n_{3844},13n_{3855},13n_{3856}\}.
\end{split}
\]

\begin{theorem}\label{thm:knots}
For every $K\in\mathcal K$ there is an epimorphism
\[
                      \pi_1(S^3\setminus K)\twoheadrightarrow Q
\]
taking meridians into $C$. Each has $\mr(K)=b(K)=4$.
\end{theorem}
\begin{proof}
The finite input retains each literal PD, its exact named-census
diagram correspondence, a four-seed propagation trace and thirteen
pairs $(p_i,v_i)$.
For each arc the checker constructs the homogeneous matrix
\eqref{eq:matrix}, checks its square is the identity and verifies
every Wirtinger conjugation
\[
                           L_oL_aL_o=L_b .
\]
Because the matrices are involutions this checks either signed
Wirtinger convention. Hence the matrices define a knot-group
homomorphism. Every $p_i$ is a nonidentity involution of $G$ and
$(I+P_{p_i})v_i=0$, so each image is in $C$.

In each case the projected four seed images generate all of $G$.
Six explicitly retained words on the four seed images evaluate to
pure translations forming a basis of $W$.
Thus every one of the twenty-four homomorphisms is surjective onto
the full $Q$, not merely a subgroup.

Theorem~\ref{thm:group} now gives $\mr(K)\ge4$.
Equivalently, direct row reduction of the four seed displacement
matrices has rank seven while each meridian displacement rank is
two, so Lemma~\ref{lem:fixed} gives the same lower bound.
The four-seed trace checks each exact diagrammatic coloring move and
reaches every arc. The published equality between bridge number and
minimum Wirtinger number \cite[Definition~1.2 and Theorem~1.3]{bkvv}
therefore gives $b(K)\le4$. Since $\mr(K)\le b(K)$, both equal four.
\end{proof}

For $13a_{373}$, the four seed arcs are $1,2,6,8$ and their images
are exactly the quartet $x_0,x_1,x_2,x_3$ above.
This provides a direct bridge between the finite-group example and
the literal knot presentation.

\section{Finite affine realization of a homological cover bound}

The same construction works for arbitrary transitive actions and finite
fields. Its purpose is to realize a homological cover bound by explicit
finite linear images; it does not improve that numerical bound.

\begin{theorem}\label{thm:realization}
Let $\Gamma$ be a classical knot group with a transitive permutation
action $P$ of degree $n>1$. Let $\mu$ be a meridian, let $c$ be the
number of cycles of $P_\mu$, and put $\delta=n-c$.
Let $Y$ be the associated branched cover, and let
$h=\dim_{\F_p}H_1(Y;\F_p)$ for a prime $p$.
There is a finite-image representation on $\F_p^{\,n+h}$ such that
every meridian has displacement rank $\delta$ and the common fixed
space has dimension one. Consequently its fixed-space bound is
\[
                \mr(K)\ge\left\lceil\frac{n-1+h}{n-c}\right\rceil.
\]
If $p\nmid n$, a direct summand of dimension $n-1+h$ has zero common
fixed space and the same meridian displacement rank.
\end{theorem}
\begin{proof}
Use a one-vertex presentation complex for the knot exterior.
Its unbranched $n$-sheeted cover has one vertex for each sheet and
one lifted edge for each generator and starting sheet.
Lift all presentation relations as two-cells.
Attach the meridian-cycle filling two-cells: for each cycle of
$P_\mu$ of length $\ell$, kill the corresponding lift of $\mu^\ell$.
These are precisely the meridian slopes filled in the branched
cover; repeated cycles on the same peripheral component give
conjugate, hence redundant, relations. Higher cells do not affect
first cohomology. Thus the resulting cellular complex computes
$H^1(Y;\F_p)$.

Choose cellular cocycles $\alpha_1,\ldots,\alpha_h$ representing a
basis of this cohomology group. Their number equals
$\dim_{\F_p}H_1(Y;\F_p)$.
For $g\in\Gamma$, define $z_j(g)\in\F_p^n$ by integrating
$\alpha_j$ along the lift of a word for $g$, with each coordinate
indexed by its ending sheet. Read words with the usual left-action
convention, so multiplication $gh$ traverses $h$ and then $g$.
Vanishing on lifted relators makes the vector independent of the
word, and path concatenation gives
\[
                     z_j(gh)=z_j(g)+P_gz_j(h).
\]
A vertex coboundary $df$ gives
$z_j(g)=(I-P_g)f$. Conversely, equality with such a principal cocycle
on the generators says exactly that all lifted-edge cochain values
are those of $df$. This proves the required identification of
cellular cohomology classes with these parabolic cocycles modulo
principal ones; no cohomological induction theorem is needed.

The meridian filling cell on each cycle says that the sum of
$z_j(\mu)$ on that cycle is zero.
For a permutation matrix over any field,
$\operatorname{im}(I-P_\mu)$ is exactly the vectors with zero sum on
each cycle: one implication telescopes, and the other follows by
successively solving the difference equation around each cycle.
No division by its length occurs.
Therefore every column $z_j(\mu)$ belongs to that image.

Put $Z_g=(z_1(g)\ \cdots\ z_h(g))$ and define
\[
                       M_g=
                       \begin{pmatrix}P_g&Z_g\\0&I_h\end{pmatrix}.
\]
The cocycle identity makes this a representation.
For a meridian, the additional columns lie in
$\operatorname{im}(I-P_\mu)$, so
$\rk(M_\mu-I)=\rk(P_\mu-I)=n-c$.

If $(v,a_1,\ldots,a_h)$ is fixed, then
\[
                    \sum_j a_jz_j(g)=(I-P_g)v
                    \qquad\text{for all }g.
\]
The corresponding combination of cellular cohomology classes is
zero. Independence of the chosen basis forces every $a_j=0$.
Now $P_gv=v$ for all $g$; transitivity makes $v$ constant.
Thus the full common fixed space is exactly the constant line.
Lemma~\ref{lem:fixed} gives the displayed lower bound.
The image is finite because the field and dimension are finite.
Also $\delta>0$: if the meridian acted trivially, all its conjugate
meridian generators would act trivially, contradicting transitivity
when $n>1$.

Finally the coordinate sum of $z_j(g)$ is a homomorphism
$\Gamma\to\F_p$, since permutations preserve that sum.
The meridian filling condition makes it zero on $\mu$, hence on
all conjugate meridian generators. Thus every $z_j(g)$ has sum zero.
When $p\nmid n$, the permutation module is the direct sum of its
constant line and its sum-zero submodule. The matrices split off
that fixed line, leaving dimension $n-1+h$ and zero common fixed
space, with the same displacement ranks.
\end{proof}

For $n=7$, $p=3$ and $h=1$, Theorem~\ref{thm:realization} yields
the seven-dimensional form used above. The explicit $Q$ certificates
add full-image and class-generating-rank information beyond the
numerical realization statement.

\section{Evidence, interpretation and limits}

The retained files are
\path{ternary_affine_quotient_inputs.json},
\path{verify_ternary_affine_quotients.py} and
\path{ternary_affine_quotient_verification.json}.
The standard-library checker verifies all twenty-four sets of
matrices and kernel words directly; it does not run a cocycle search
or use a cover-group certificate. Each knot has its own six kernel
words proving the full image. One independent affine enumeration
checks the order $122472$, class size $189$ and two-generator claim.
Corrupting a matrix, a kernel word or a named-census map is rejected;
the zero-translation principal lift is also rejected because its
common fixed space is nonzero.

The named diagrams are bound to the retained independently packaged
census correspondence and its actual census-read receipt.
That identity dependency and the published geometric
bridge--Wirtinger theorem remain separate from the new finite-group
calculation. No stale or missing cover-group certificate is a
premise of the matrix proof.

The lifts were found from a nonprincipal parabolic cocycle in the
seven-point permutation module over $\F_3$.
All involutory translations have coordinate sum zero, so the
six-dimensional module $W$ suffices.
Their discovery through seven-sheeted homology is not needed for
verification: the matrix equations and pure translation words are
complete finite witnesses. No claim of a new general cocycle
principle or independent worldwide priority is made.

The result separates ordinary group generation from generation by
the distinguished meridian class. It supplies the same sharp lower
bound four on these knots; it does not assert connected-sum
additivity of this finite-class bound, a supremum over all finite
groups, or a counterexample to MRC.
The imported geometric theorem is not reproved, and these exact
internal replays are not formal verification or external specialist
review.

\begin{thebibliography}{9}
\bibitem{bkvv}
R. Blair, A. Kjuchukova, R. Velazquez and P. Villanueva,
Wirtinger systems of generators of knot groups,
Commun. Anal. Geom. 28 (2020), 243--262,
Definition~1.2 and Theorem~1.3, printed pp.~245--246.
\url{https://archive.intlpress.com/site/pub/files/_fulltext/journals/cag/2020/0028/0002/CAG-2020-0028-0002-a002.pdf}.
\end{thebibliography}
\end{document}
